\emph{Planar model.} \texttt{Human0914} moves in the sagittal plane only, with seven muscles per leg, no subtalar or metatarsophalangeal joint and no arms. Patients with drop foot often clear the foot by hip circumduction or hip hiking; the model cannot, so it shows only the sagittal compensations (steppage and changed timing), which exaggerates them, as Dermksian et al.\ also noted \cite{dermksian}. Without a toe joint, toe clearance is the height of a point fixed to the foot.

\emph{A reflex controller is not a patient.} Drop foot here is a weaker TA with intact sensing and reflexes; a peroneal nerve lesion or a stroke also changes sensation, spasticity and the other muscles. ``Adaptation'' is a re-optimization of reflex parameters for effort and speed, run to convergence. A person adapts for other goals (stability, fear of falling, comfort), partly within a session, and does not reach an optimum. The frozen and re-optimized conditions bracket the first day and a fully adapted user; neither is a clinical prediction. The optimized gaits also have little reserve: the healthy controller falls with a 2~N\,m pulse in swing, and the adapted gaits fall after modest pushes, which is not how people walk.

\emph{Idealized attachment.} The device applies an exact moment pair to the shank and the foot. There is no strap or cuff compliance, no misalignment between the device and the anatomical joint, no added mass or inertia, and no friction. Real orthoses lose part of their torque to soft tissue and shift during use; the added distal mass alone would raise the swing cost.

\emph{Optimizer randomness.} CMA-ES finds different local optima from different seeds: the 50\% runs differ in cost by 40\% between seeds. Three seeds show this spread but do not measure it well, and an SD over three values is a weak basis for the inconclusive band of the decision rule. Each run was stopped at a fixed number of generations, and convergence was checked on the cost history only.

\emph{Device chain.} The actuator was sized and its controller designed offline, against the normative ankle trajectory rather than the model's own, and the SCONE device sees only the fitted delay and lag, not the full SEA dynamics, motor heating or the drive. The detector was tuned on simulated gait of a model whose strides are nearly identical; the real-data test shows larger, subject-dependent errors. The active controller's gains came from a sizing calculation and only the swing target was chosen, by a rule fixed in advance; other gains would give other results.

\emph{Measures.} Ten-second simulations give four analysed strides per run. The metabolic cost is the Wang et al.\ model \cite{wang2012} inside the objective, not a measured quantity, and the device's electrical energy is reported separately and not added to it.
